% !TeX program = xelatex
% Reusable source template, not a completed paper summary.
% Replace all instructional prose and source metadata before delivery.
\documentclass[UTF8,fontset=fandol,zihao=-4,a4paper]{ctexart}
\usepackage[margin=25mm]{geometry}
\usepackage{fontspec}
% Load TeX-distributed font files without requiring OS font registration.
\setmainfont{texgyretermes-regular.otf}[
  BoldFont = texgyretermes-bold.otf,
  ItalicFont = texgyretermes-italic.otf,
  BoldItalicFont = texgyretermes-bolditalic.otf
]
\usepackage{amsmath,amssymb,bm}
\usepackage{booktabs,tabularx,graphicx}
\usepackage{setspace}
\usepackage[hidelinks,unicode]{hyperref}
\urlstyle{same}

\onehalfspacing
\setlength{\parindent}{2\ccwd}
\setlength{\parskip}{0pt}
\clubpenalty=10000
\widowpenalty=10000
\pagestyle{plain}
\ctexset{
  section = {
    number = \arabic{section},
    format = \raggedright\zihao{4}\bfseries,
    beforeskip = 12bp,
    afterskip = 6bp
  },
  subsection = {
    number = \arabic{section}.\arabic{subsection},
    format = \raggedright\zihao{-4}\bfseries,
    beforeskip = 6bp,
    afterskip = 3bp
  }
}
\DeclareMathOperator{\argmin}{arg\,min}
\DeclareMathOperator{\argmax}{arg\,max}

\begin{document}

\begin{center}
  {\zihao{2}\bfseries 文献中文题名：阅读总结\par}
  \vspace{6bp}
  {\zihao{5} 文献阅读总结\par}
\end{center}

\noindent\textbf{原文标题：}填写核验后的原文完整标题。\par
\noindent\textbf{原文作者与发表信息：}填写作者、发表来源及年份。\par
\noindent\textbf{原文链接或 DOI：}填写已核验的链接或 DOI。\par
\noindent\textbf{阅读范围：}注明全文、摘要或指定章节，以及页码规则。\par

\begin{abstract}
用中文连贯概括原文研究问题、方法、核心发现及局限。仅使用已核验信息，不把本总结表述为独立完成的原创研究。
\end{abstract}
\noindent\textbf{关键词：}填写三至五个中文关键词，以分号分隔。\par

\section{引言}
说明研究背景、原文要解决的问题和阅读本总结所需的基础概念，并引用原文\cite{source}。

\section{核心方法或理论框架}
按适合原文的逻辑解释方法、理论假设或综述组织方式。需要二级标题时使用 subsection，不手工编号。

% Add equations only when supported by the source. Example syntax:
% \begin{equation}
%   \hat{\theta} = \argmin_{\theta} \mathcal{L}(\theta)
%   \label{eq:objective}
% \end{equation}
% Explain every variable, unit, assumption, and the original equation locator.
% Refer to it with \eqref{eq:objective}; preserve the source's actual symbols.

\section{主要结果或核心论证}
依据原文证据解释主要发现，保留必要的数值、单位、实验条件和比较基线。理论或综述论文改用对应的命题、论证或研究进展。

% Optional table syntax, to be filled only with verified data:
% \begin{table}[htbp]
%   \centering
%   \caption{Describe the comparison, units, conditions, and source in Chinese.}
%   \label{tab:comparison}
%   \begin{tabularx}{\linewidth}{lXX}
%     \toprule
%     Method & Metric & Conditions \\
%     \midrule
%     ... & ... & ... \\
%     \bottomrule
%   \end{tabularx}
% \end{table}

\section{讨论与局限}
区分原文明确报告的局限和总结者的解释，注明证据与适用范围；缺失的信息如实说明。

\section{结论}
用中文归纳原文能够支持的结论、学习收获及仍未解决的问题，不增加新证据或夸大贡献。

\begin{thebibliography}{99}
\bibitem{source}
原文作者．原文题名．发表来源，年份．填写核验后的 DOI 或来源链接。
% Add only sources actually consulted and verified, using consecutive bibitems.
\end{thebibliography}

\end{document}
